% GENERATED FILE. Edit canonical lessons, then run npm run generate:books.
% canonical-source-hash: fnv1a64:d8c750f1c354c76c
% canonical-lessons: RU-C396-dorozhat, RU-C396-deshevet, RU-C396-inflyaciya, RU-C396-kurs, RU-C396-obmen

\chapter{Rising prices}
\label{ch:ru-rising-prices}
\hlchaptermodality{396}

\begin{chapteropening}
\textbf{By the end of this chapter:} \emph{I can say that prices go up or down, talk about inflation and the exchange rate, and change money.}

Complete the last lesson of chapter 396: I can say that prices go up or down, talk about inflation and the exchange rate, and change money.
\end{chapteropening}

\section[\texorpdfstring{\ru{дорожать} (dorozhát)}{дорожать (dorozhát)}]{\ru{дорожать} (dorozhát) — to go up in price}

\label{lesson:RU-C396-dorozhat}



\begin{quote}
Before the new one: say the Russian for cost, then the Russian for a budget.
\end{quote}

\subsection*{You'll want to know: \ru{дорожать}}
\textbf{\ru{дорожать}} (\emph{dorozhát}) — \textquotedblleft{}to go up in price\textquotedblright{}.

\begin{quote}
\ru{Хлеб} \ru{опять} \ru{дорожает}.
\end{quote}

\begin{quote}
\emph{Khleb opyát dorozháyet.}
\end{quote}

\begin{quote}
Bread is going up in price again.
\end{quote}

Imperfective; perfective \emph{\ru{подорожать}}: \emph{\ru{всё} \ru{подорожало}} — everything has gone up.

\begin{cousinweb}
\textbf{\ru{дорожать}} is built on \emph{\ru{дорогой}} \textquotedblleft{}dear, expensive\textquotedblright{}, inherited from Proto-Slavic \emph{*dorgŭ} \textquotedblleft{}dear\textquotedblright{}: to become dear.
\end{cousinweb}

\begin{grammarlens}[title={What you've built}]
One more word for telling, explaining and describing at B1.
\end{grammarlens}

\subsection*{Your turn}
\begin{itemize}
\raggedright
  \item \emph{Say it:} \emph{dorozhát}
  \item \emph{Say it:} \emph{dorozhát}, once more
  \item \emph{Say it:} a sentence of your own with \emph{dorozhát}
  \item \emph{Recall:} say the Russian for cost, then the Russian for a budget, then say \emph{dorozhát} again
\end{itemize}

\begin{tcolorbox}[breakable,colback=teal!4,colframe=teal!35!black,title={Before you move on}]
What does \ru{дорожать} mean? (\textquotedblleft{}To go up in price\textquotedblright{}.) Say it once more.
\end{tcolorbox}
\section[\texorpdfstring{\ru{дешеветь} (deshevét)}{дешеветь (deshevét)}]{\ru{дешеветь} (deshevét) — to get cheaper}

\label{lesson:RU-C396-deshevet}



\begin{quote}
Before the new one: say the Russian for a budget, then the Russian for to go up in price.
\end{quote}

\subsection*{You'll want to know: \ru{дешеветь}}
\textbf{\ru{дешеветь}} (\emph{deshevét}) — \textquotedblleft{}to get cheaper\textquotedblright{}.

\begin{quote}
\ru{Летом} \ru{фрукты} \ru{дешевеют}.
\end{quote}

\begin{quote}
\emph{Létom frúkty deshevéyut.}
\end{quote}

\begin{quote}
In summer fruit gets cheaper.
\end{quote}

Imperfective; perfective \emph{\ru{подешеветь}}.

\begin{cousinweb}
\textbf{\ru{дешеветь}} is built on \emph{\ru{дешёвый}} \textquotedblleft{}cheap\textquotedblright{}: to become cheap. \emph{\ru{Дешёвый}} goes back to Old Russian \emph{deševŭ}, which is usually derived from \emph{\ru{десить}} \textquotedblleft{}to find, to come upon\textquotedblright{}: what is cheap is easily come by.
\end{cousinweb}

\begin{grammarlens}[title={What you've built}]
One more word for telling, explaining and describing at B1.
\end{grammarlens}

\subsection*{Your turn}
\begin{itemize}
\raggedright
  \item \emph{Say it:} \emph{deshevét}
  \item \emph{Say it:} \emph{deshevét}, once more
  \item \emph{Say it:} a sentence of your own with \emph{deshevét}
  \item \emph{Recall:} say the Russian for a budget, then the Russian for to go up in price, then say \emph{deshevét} again
\end{itemize}

\begin{tcolorbox}[breakable,colback=teal!4,colframe=teal!35!black,title={Before you move on}]
What does \ru{дешеветь} mean? (\textquotedblleft{}To get cheaper\textquotedblright{}.) Say it once more.
\end{tcolorbox}
\section[\texorpdfstring{\ru{инфляция} (inflyátsiya)}{инфляция (inflyátsiya)}]{\ru{инфляция} (inflyátsiya) — inflation}

\label{lesson:RU-C396-inflyaciya}



\begin{quote}
Before the new one: say the Russian for to go up in price, then the Russian for to get cheaper.
\end{quote}

\subsection*{You'll want to know: \ru{инфляция}}
\textbf{\ru{инфляция}} (\emph{inflyátsiya}) — \textquotedblleft{}inflation\textquotedblright{}.

\begin{quote}
\ru{В} \ru{прошлом} \ru{году} \ru{инфляция} \ru{была} \ru{высокой}.
\end{quote}

\begin{quote}
\emph{V próshlom godú inflyátsiya bylá vysókoy.}
\end{quote}

\begin{quote}
Last year inflation was high.
\end{quote}

Feminine. \emph{\ru{Уровень} \ru{инфляции}} — the rate of inflation.

\begin{cousinweb}
\textbf{\ru{инфляция}} was borrowed through western languages from Latin \emph{īnflātiō} \textquotedblleft{}a blowing up, a swelling\textquotedblright{} (EtymDB; Lewis \& Short), from \emph{īnflāre} \textquotedblleft{}to blow into\textquotedblright{}: prices swell.
\end{cousinweb}

\begin{grammarlens}[title={What you've built}]
One more word for telling, explaining and describing at B1.
\end{grammarlens}

\subsection*{Your turn}
\begin{itemize}
\raggedright
  \item \emph{Say it:} \emph{inflyátsiya}
  \item \emph{Say it:} \emph{inflyátsiya}, once more
  \item \emph{Say it:} a sentence of your own with \emph{inflyátsiya}
  \item \emph{Recall:} say the Russian for to go up in price, then the Russian for to get cheaper, then say \emph{inflyátsiya} again
\end{itemize}

\begin{tcolorbox}[breakable,colback=teal!4,colframe=teal!35!black,title={Before you move on}]
What does \ru{инфляция} mean? (\textquotedblleft{}Inflation\textquotedblright{}.) Say it once more.
\end{tcolorbox}
\section[\texorpdfstring{\ru{курс} (kurs)}{курс (kurs)}]{\ru{курс} (kurs) — an exchange rate; a course; a year (of study)}

\label{lesson:RU-C396-kurs}



\begin{quote}
Before the new one: say the Russian for to get cheaper, then the Russian for inflation.
\end{quote}

\subsection*{You'll want to know: \ru{курс}}
\textbf{\ru{курс}} (\emph{kurs}) — \textquotedblleft{}an exchange rate; a course; a year (of study)\textquotedblright{}.

\begin{quote}
\ru{Какой} \ru{сегодня} \ru{курс} \ru{доллара}?
\end{quote}

\begin{quote}
\emph{Kakóy segódnya kurs dóllara?}
\end{quote}

\begin{quote}
What is the dollar exchange rate today?
\end{quote}

Masculine. \emph{\ru{На} \ru{третьем} \ru{курсе}} — in the third year at university. \emph{\ru{Курс} \ru{английского}} — an English course.

\begin{cousinweb}
\textbf{\ru{курс}} is usually traced to German \emph{Kurs}, from Latin \emph{cursus} \textquotedblleft{}a running, a course\textquotedblright{} (Lewis \& Short), from \emph{currere} \textquotedblleft{}to run\textquotedblright{}.
\end{cousinweb}

\begin{grammarlens}[title={What you've built}]
One more word for telling, explaining and describing at B1.
\end{grammarlens}

\subsection*{Your turn}
\begin{itemize}
\raggedright
  \item \emph{Say it:} \emph{kurs}
  \item \emph{Say it:} \emph{kurs}, once more
  \item \emph{Say it:} a sentence of your own with \emph{kurs}
  \item \emph{Recall:} say the Russian for to get cheaper, then the Russian for inflation, then say \emph{kurs} again
\end{itemize}

\begin{tcolorbox}[breakable,colback=teal!4,colframe=teal!35!black,title={Before you move on}]
What does \ru{курс} mean? (\textquotedblleft{}An exchange rate; a course; a year (of study)\textquotedblright{}.) Say it once more.
\end{tcolorbox}
\section[\texorpdfstring{\ru{обмен} (obmén)}{обмен (obmén)}]{\ru{обмен} (obmén) — an exchange}

\label{lesson:RU-C396-obmen}



\begin{quote}
Before the new one: say the Russian for inflation, then the Russian for an exchange rate.
\end{quote}

\subsection*{You'll want to know: \ru{обмен}}
\textbf{\ru{обмен}} (\emph{obmén}) — \textquotedblleft{}an exchange\textquotedblright{}.

\begin{quote}
\ru{Пункт} \ru{обмена} \ru{работает} \ru{до} \ru{восьми}.
\end{quote}

\begin{quote}
\emph{Punkt obména rabótayet do vosmí.}
\end{quote}

\begin{quote}
The exchange office is open until eight.
\end{quote}

Masculine. \emph{\ru{Обмен} \ru{опытом}} — sharing experience.

\begin{cousinweb}
\textbf{\ru{обмен}} is formed from \emph{\ru{обменять}} \textquotedblleft{}to exchange\textquotedblright{}, from \emph{\ru{менять}} \textquotedblleft{}to change\textquotedblright{}, inherited from Proto-Slavic \emph{*měniti}.
\end{cousinweb}

\begin{grammarlens}[title={What you've built}]
One more word for telling, explaining and describing at B1.
\end{grammarlens}

\subsection*{Your turn}
\begin{itemize}
\raggedright
  \item \emph{Say it:} \emph{obmén}
  \item \emph{Say it:} \emph{obmén}, once more
  \item \emph{Say it:} a sentence of your own with \emph{obmén}
  \item \emph{Recall:} say the Russian for inflation, then the Russian for an exchange rate, then say \emph{obmén} again
\end{itemize}

\begin{tcolorbox}[breakable,colback=teal!4,colframe=teal!35!black,title={Before you move on}]
What does \ru{обмен} mean? (\textquotedblleft{}An exchange\textquotedblright{}.) Say it once more.
\end{tcolorbox}
